% ============================================================
% Feature Split Evaluation Table — Information Gain
% CSC491 Textbook, Chapter: Machine Learning
%
% Purpose: Illustrates how a decision tree evaluates each feature
% as a candidate split, connecting to the hieroglyph dataset in
% tab:hieroglyph-features. Shows resulting group purity to motivate
% Information Gain as a formal measure of the best split.
%
% Preamble requirements:
%     \usepackage{booktabs, array, xcolor, multirow, fontspec}
%     \newfontfamily\egyptfont{Noto Sans Egyptian Hieroglyphs}
%     \definecolor{puregreen}{RGB}{0, 120, 80}
%     \definecolor{impurered}{RGB}{180, 40, 40}
%     \definecolor{midgold}{RGB}{160, 110, 0}
% ============================================================

\begin{table}[ht]
\centering
\setlength{\tabcolsep}{8pt}
\renewcommand{\arraystretch}{1.9}
\caption{%
  Evaluating candidate splits on the hieroglyph dataset
  (see Table~\ref{tab:hieroglyph-features}).
  For each feature, the glyphs are divided into two groups based on
  the split condition. A good split produces groups where each
  contains mostly one class --- the tree prefers the feature that
  results in the purest groups, quantified by \textbf{Information Gain}.%
}
\label{tab:infogain-splits}
\begin{tabular}{
  >{\raggedright\arraybackslash}p{2.4cm}   % feature
  >{\centering\arraybackslash}p{1.8cm}     % split condition
  >{\centering\arraybackslash}p{3.0cm}     % group: yes
  >{\centering\arraybackslash}p{3.0cm}     % group: no
  >{\centering\arraybackslash}p{2.2cm}     % purity verdict
}
\toprule
\textbf{Feature}
  & \textbf{Split}
  & \textbf{Group A} (\textit{Yes})
  & \textbf{Group B} (\textit{No})
  & \textbf{Purity} \\
\midrule

% --- Animal? ---
\textbf{Animal?}
  & Animal = Yes
  & {\egyptfont 𓅓} Owl \quad {\egyptfont 𓆑} Viper \quad {\egyptfont 𓃭} Lion
  & {\egyptfont 𓂀} Eye \quad {\egyptfont 𓈖} Water \quad {\egyptfont 𓏏} Bread
    \quad {\egyptfont 𓇋} Reed \quad {\egyptfont 𓎛} Rope
  & \textcolor{puregreen}{\textbf{High}} \\

\addlinespace[2pt]
\midrule
\addlinespace[2pt]

% --- Has Eye? ---
\textbf{Has Eye?}
  & Has Eye = Yes
  & {\egyptfont 𓂀} Eye \quad {\egyptfont 𓅓} Owl \quad {\egyptfont 𓃭} Lion
  & {\egyptfont 𓈖} Water \quad {\egyptfont 𓆑} Viper \quad {\egyptfont 𓏏} Bread
    \quad {\egyptfont 𓇋} Reed \quad {\egyptfont 𓎛} Rope
  & \textcolor{midgold}{\textbf{Medium}} \\

\addlinespace[2pt]
\midrule
\addlinespace[2pt]

% --- Vertical? ---
\textbf{Vertical?}
  & Vertical = Yes
  & {\egyptfont 𓇋} Reed
  & {\egyptfont 𓂀} Eye \quad {\egyptfont 𓅓} Owl \quad {\egyptfont 𓈖} Water
    \quad {\egyptfont 𓆑} Viper \quad {\egyptfont 𓏏} Bread
    \quad {\egyptfont 𓃭} Lion \quad {\egyptfont 𓎛} Rope
  & \textcolor{midgold}{\textbf{Medium}} \\

\addlinespace[2pt]
\midrule
\addlinespace[2pt]

% --- Curved Strokes (continuous, split at median) ---
\textbf{Curved} $\geq 3$
  & Strokes $\geq 3$
  & {\egyptfont 𓂀} Eye \quad {\egyptfont 𓅓} Owl \quad {\egyptfont 𓈖} Water
    \quad {\egyptfont 𓆑} Viper \quad {\egyptfont 𓃭} Lion
  & {\egyptfont 𓏏} Bread \quad {\egyptfont 𓇋} Reed \quad {\egyptfont 𓎛} Rope
  & \textcolor{impurered}{\textbf{Low}} \\

\bottomrule
\end{tabular}

\smallskip
\footnotesize
\textit{Note:} ``Purity'' here is qualitative --- in the next section we will replace this
judgment with a precise number using \textbf{Gini Impurity} and \textbf{Information Gain}.
The \textit{Animal?} feature produces the cleanest separation: Group~A contains exactly
the three animal glyphs, and Group~B contains everything else.
A decision tree formalizes this intuition by selecting the feature with the highest
information gain --- the split that reduces entropy the most.

\end{table}